\chapter{Sistemi di Particelle Identiche}

	Si consideri un sistema composto del tipo $\mathcal{H}_1 \otimes \mathcal{H}_2$. Esso ha dimensione:
	
\begin{align*}
	\dim \mathcal{H} = n_1 n_2
\end{align*}

	che sono rispettivamente la dimensione di $\mathcal{H}_1$ e $\mathcal{H}_2$.
	\\	
	\\
	Mi metto nel caso in cui vale: $\mathcal{H}_1 = \mathcal{H}_2$, cioè nel caso in cui gli autospazi sono descritti dalle osservabili. Se ho una soluzione per $\mathcal{H}$, posso allora scambiare i pedici $1$ e $2$ ed avere di nuovo una soluzione per $\mathcal{H}$: si tratta di una \textit{simmetria di scambio.}
	\\
	\\
	In meccanica classica, date $2$ particelle ad esempio, se fisso le condizioni iniziali potrò sempre distinguere l'una dall'altra siccome soddisfano una legge del moto definita e seguono quindi delle traiettorie fisse nello spazio delle fasi. Invece, in meccanica quantistica, esistono \textbf{sistemi di particelle identiche}, ossia \textit{indistinguibili} dal momento che non esistono traiettorie. Non ho dunque modo di distinguere fisicamente la situazione $\qty(1, 2)$ o $\qty(2, 1)$. Si tratta del \textbf{principio di indistinguibilità}. Un esempio di sistema di particelle identiche è un sistema di due elettroni.

\section{Principio di Pauli: bosoni e fermioni}

	Generalizzando ad $N$ particelle, si abbiano $N$ autospazi di dimensione $k$, ovvero tali che:

\begin{align*}	
	\dim \mathcal{H}_k = k \, \, \Rightarrow \, \, \dim \mathcal{H}_N = \dim \qty(\mathcal{H}_k \otimes \ldots \otimes \mathcal{H}_k)
\end{align*}

	se siamo nel caso in cui tali sottospazi sono \textit{distinguibili}, allora:
	
\begin{align*}
	\dim \mathcal{H}_k = k^N
\end{align*}

	altrimenti devo togliere tutti i sottospazi identici ad altri. In altre parole, se una base di $\mathcal{H}_k$ è:
	
\begin{align*}
	\qty{ \ket{\alpha}, \, \alpha = 1, \ldots, k}
\end{align*}

	una base di $\mathcal{H}_N$ sarà:
	
\begin{align*}
	\qty{\ket{\alpha_1}, \ldots, \ket{\alpha_N}, \, \alpha_i = 1, \ldots, k}
\end{align*}
	
	quindi ad esempio se non distinguo $\ket{\alpha_1}\ket{\alpha_2}$ da $\ket{\alpha_2} \ket{\alpha_1}$, dovrò tenere conto della presenza di questi due stati indistinguibili. Introduco allora il seguente operatore:
	
\begin{mydef}	
	Definiamo l'\textbf{operatore di scambio} $\hat{P}$, ossia una matrice le cui entrate sono:
	
\begin{align*}
	\hat{P}_{ij} : \mathcal{H}_N \to \mathcal{H}_N
\end{align*}

	tale che:
	
\begin{align*}
	\ket{\alpha_1} \cdots \ket{\alpha_i} \cdots \ket{\alpha_j} \cdots \ket{\alpha_N} \to \ket{\alpha_1} \cdots \ket{\alpha_j} \cdots \ket{\alpha_i} \cdots \ket{\alpha_N}
\end{align*}

	e che è unitaria. Pertanto vale:

\begin{align*}
	\hat{P}_{ij}^2 = \mathbb{I} \qquad ; \qquad \hat{P}_{ij}^{\dagger} = \hat{P}_{ij}
\end{align*}

\end{mydef}

	Enunciamo ora il seguente ed importantissimo principio, il \textbf{principio di Pauli}:
	
\begin{myprinc}

	Il \textbf{principio di Pauli} si articola nei seguenti $3$ punti:

\begin{enumerate}[label=\arabic*)]

	\item In un sistema di particelle identiche, le \textbf{osservabili fisiche} sono una sottoalgebra del sistema complessivo: tutte le \textit{osservabili fisiche} sono quelle in cui operatori associati commutano con $\hat{P}_{ij}$:
	\begin{align*}
		\hat{O} : \mathcal{H}_N \to \mathcal{H}_N \, \, \Rightarrow \, \, \qty[\hat{O}, \hat{P}_{ij}] = 0
	\end{align*}

	\item Come agisce $\hat{P}_{ij}$ sullo stato $\psi$? 
	
		Ho che:
	
	\begin{align*}
		\hat{P}_{ij} \psi = \lambda \psi \, \, \Rightarrow \, \, \hat{P}_{ij}^2 \psi = \lambda \hat{P}_{ij} \psi \, \, \Rightarrow \, \, \hat{P}^2_{ij} \psi = \lambda^2 \psi 
	\end{align*}
	
		ed essendo che $\hat{P}^2_{ij} = \mathbb{I}$ per ipotesi allora:
		
	\begin{align*}
		\psi = \lambda^2 \psi  \, \, \Rightarrow \, \, \lambda = \pm 1
	\end{align*}

		pertanto, poiché dal primo punto so che $\mathcal{H}_{\text{fisici}} \subset \mathcal{H}$, in particolare ho che la sottoalgebra è composta da stati simmetrici e stati antisimmetrici:
		
	\begin{align*}
		& \mathcal{H}_+ = \qty{\psi \in \mathcal{H}_N, \, \hat{P}_{ij}\psi = \psi, \, \forall i, j } \, \, \Rightarrow \, \, \textbf{simmetrici} \\
		& \mathcal{H}_- = \qty{\psi \in \mathcal{H}_N, \, \hat{P}_{ij}\psi = - \psi, \, \forall i, j } \, \, \Rightarrow \, \, \textbf{antisimmetrici}
	\end{align*}
	
	\item Quando $\mathcal{H}_k$ è lo spazio degli stati di una particella non relativistica che si muove in uno spazio tridimensionale con spin $s$, cioè quando:
	
	\begin{align*}
		\mathcal{H}_k = L_2\qty(\mathbb{R}^3) \otimes \mathcal{H}_s	
	\end{align*}
		
		si ha allora:
	
	\begin{itemize}
		\item $\mathcal{H}_k = \mathcal{H}_{+}$ se $s$ è \textbf{intero}: in questo caso la particella segue la \textit{statistica di Bose-Einstein}, per questo si parla di \textbf{bosoni}. La dimensione dello spazio è:
		\begin{align*}
			\dim \mathcal{H}_+ = \mqty(N + k -1 \\ k -1)
		\end{align*}
		
		\item $\mathcal{H}_k = \mathcal{H}_{-}$ se $s$ è \textbf{semintero}: in questo caso invece la particella segue la \textit{statistica di Fermi-Dirac}, per questo si parla di \textbf{fermioni}. La dimensione dello spazio è:
		\begin{align*}
			\dim \mathcal{H}_- = \mqty(k \\ N) \quad \text{se} \, \, \, N \leq k
		\end{align*}
			altrimenti se $ N \geq k$, la dimensione è nulla.
	\end{itemize}
	
	tutto ciò va sotto al nome: \textbf{Connessione Spin - Statistica} che in meccanica quantistica non relativistica è un assioma, mentre in meccanica quantistica relativistica è un teorema vero e proprio.

\end{enumerate}
\end{myprinc}

	Consideriamo ora anche lo spin. 
	\\
	\\
	Ad esempio, per $N = 2$, nel caso dei bosoni si ha che:
	
\begin{align*}
	\Psi_+ \qty(\vb{r}_1, \vb{r}_2, \sigma_1, \sigma_2) = \frac{1}{\sqrt{2}} \, \qty\Big( \psi_a\qty(\vb{r}_1, \sigma_1) \, \psi_b\qty(\vb{r}_2, \sigma_2) + \,  \psi_a\qty(\vb{r}_2, \sigma_2) \, \psi_b \qty(\vb{r}_1, \sigma_1)     )
\end{align*}

	mentre nel caso di fermioni:

\begin{align*}
	\Psi_- \qty(\vb{r}_1, \vb{r}_2, \sigma_1, \sigma_2) = \frac{1}{\sqrt{2}} \, \qty\Big( \psi_a\qty(\vb{r}_1, \sigma_1) \, \psi_b\qty(\vb{r}_2, \sigma_2) - \,  \psi_a\qty(\vb{r}_2, \sigma_2) \, \psi_b \qty(\vb{r}_1, \sigma_1)     )
\end{align*}

	dove $a$ e $b$ etichettano le due particelle: dal principio di indistinguibilità non posso sapere chi è uno e chi sia l'altro, pertanto la funzione d'onda sarà una combinazione lineare dei due possibili casi (particella $a$ in $\qty(\vb{r}_1, \sigma_1)$ e particella $b$ in $\qty(\vb{r}_2, \sigma_2)$ e viceversa), la cui simmetria sarà dettata dall'algebra del sottospazio fisico a cui appartengono (pari per i bosoni e dispari per i fermioni). 
	\\
	\\
	Da ciò deriva il più noto, anche se è conseguenza diretta del principio generale, \textbf{principio di esclusione}:
	
\begin{myprinc}
	\textit{Due fermioni identici non possono occupare simultaneamente lo stesso stato.}
\end{myprinc}

	La simmetria della funzione d'onda è dettata dall'algebra del sottospazio a cui le particelle appartengono: per i fermioni, la funzione d'onda deve essere dispari (antisimmetrica) rispetto allo scambio di due particelle. Questo significa che se scambi due fermioni identici, la funzione d'onda cambia segno:
	
\begin{align*}
	\Psi\qty(2, 1) = - \Psi\qty(1, 2)
\end{align*}

	Per soddisfare questa proprietà antisimmetrica, la funzione d'onda totale $\Psi_-$ si costruisce sottraendo le due possibili configurazioni:
	
\begin{align*}
	\Psi_- \sim \psi_a\qty(1)\psi_b\qty(2) - \psi_a\qty(2)\psi_b\qty(1)
\end{align*}

	 Se facessimo finta che non valga il \textit{principio di Pauli}, potremmo avere che nell'esempio di prima $\Psi_-\qty(\vb{r}_1, \vb{r}_2) = 0$, poiché ho il termine $\psi_a\qty(\vb{r}_1, \sigma_1) \psi_b\qty(\vb{r}_1, \sigma_1)$ meno se stesso, che è impossibile. Infatti se la funzione d'onda $\Psi$ è identicamente zero, significa che la probabilità di trovare il sistema in quella configurazione è nulla (la particella non esiste). Quindi, fisicamente, quello stato non può realizzarsi.
	\\
	\\
	In generale, per $n$ particelle identiche, posso fattorizzare la funzione d'onda in una parte spaziale ed una di spin (\textbf{spinore}):
	
\begin{align*}
	\Psi\qty(\vb{r}_1, \ldots, \vb{r}_n, \sigma_1, \ldots, \sigma_n) = \phi\qty(\vb{r}_1, \ldots, \vb{r}_n) \, \chi\qty(\sigma_1, \ldots, \sigma_n)
\end{align*}

	Prendiamo per esempio il caso di $n  = 2$ fermioni identici. Poiché l'operatore di scambio agisce simultaneamente sulle due funzioni, per garantire l'antisimmetria complessiva della funzione, ho che:

\begin{itemize}
	\item Se la parte spaziale è simmetrica, la parte di spin deve essere antisimmetrica allora:
	\begin{align*}
		\psi\qty(\vb{r}_1, \vb{r}_2) = \psi\qty(\vb{r}_2, \vb{r}_1) \qquad ; \qquad \chi\qty(\sigma_1, \sigma_2) = - \chi\qty(\sigma_2, \sigma_1)
	\end{align*}

	\item Se la parte spaziale invece è antisimmetrica, la parte di spin deve essere simmetrica:
	\begin{align*}
		\psi\qty(\vb{r}_1, \vb{r}_2) = - \psi\qty(\vb{r}_2, \vb{r}_1) \qquad ; \qquad \chi\qty(\sigma_1, \sigma_2) = \chi\qty(\sigma_2, \sigma_1)
	\end{align*}
	
\end{itemize}
		
	Discorso analogo vale per $n = 2$ bosoni identici. Poiché l'operatore di scambio agisce simultaneamente sulle due funzioni, per garantire la simmetria complessiva della funzione ho che:
	
\begin{itemize}
	
	\item Se la parte spaziale è simmetrica, la parte di spin deve essere simmetrica:
	\begin{align*}
		\psi\qty(\vb{r}_1, \vb{r}_2) =  \psi\qty(\vb{r}_2, \vb{r}_1) \qquad ; \qquad \chi\qty(\sigma_1, \sigma_2) = \chi\qty(\sigma_2, \sigma_1)
	\end{align*}
	
	\item Se la parte spaziale è antisimmetrica, la parte di spin deve essere antisimmetrica:
	\begin{align*}
		\psi\qty(\vb{r}_1, \vb{r}_2) =  -\psi\qty(\vb{r}_2, \vb{r}_1) \qquad ; \qquad \chi\qty(\sigma_1, \sigma_2) = -\chi\qty(\sigma_2, \sigma_1)
	\end{align*}
	
\end{itemize}

	Prima di andare avanti, è opportuno fare alcune considerazioni molto delicate, che avranno però importanti conseguenze. 
	\\
	\\
	Consideriamo due particelle in moto in un potenziale centrale, che dipende soltanto dalla distanza relativa tra di esse. L'Hamiltoniana che descrive il sistema è:
	
\begin{align*}
	\hat{H} = \frac{\hat{\vb{p}}^2_1}{2m_1} + \frac{\hat{\vb{p}}^2_2}{2m_2} + V\qty(r)
\end{align*}

	dove $r = |\vb{r}_1 - \vb{r}_2|$ è la coordinata relativa tra le due particelle. Definisco ora la coordinata del \textit{centro di massa} come:
	
\begin{align*}
	\vb{R}_{\text{CM}} = \frac{m_1 \vb{r}_1 + m_2 \vb{r}_2}{m_1 + m_2}
\end{align*}

	da cui ho:
	
\begin{align*}
	\vb{V}_{\text{CM}} = \frac{m_1 \vb{v}_1 + m_2 \vb{v}_2}{m_1 + m_2}
\end{align*}

	e allora posso scrivere la quantità di moto rispetto al centro di massa come:
	
\begin{align*}
	\vb{P}_{\text{CM}} = \qty(m_1 + m_2) \, \vb{V}_{\text{CM}} = \vb{p}_1 + \vb{p}_2
\end{align*}

	Ricordiamo ora del problema a due corpi, le singole quantità di moto sono espresse in termini di quantità di moto relativa, ovvero:
	
\begin{align*}
	\vb{p}_{\text{rel}} = \mu \, \vb{V}_{\text{rel}}
\end{align*}

	dove $\mu$ è la \textit{massa ridotta} e allora possiamo scrivere:

\begin{align*}
	\vb{p}_{\text{rel}} = \mu \, \vb{V}_{\text{rel}} = \frac{m_1 m_2}{m_1 + m_2} \, \vb{V}_{\text{rel}} = \frac{m_1 m_2}{m_1 + m_2} \, |\vb{v}_1 - \vb{v}_2| = \frac{m_2 \vb{p}_1 - m_1 \vb{p}_2}{ m_1 + m_2 } 
\end{align*}

	dove infatti:
	
\begin{align*}
	& \vb{p}_1 = m_1 \vb{v}_1 \, \Rightarrow \, \vb{v}_1 = \frac{\vb{p}_1}{m_1} \\
	& \vb{p}_2 = m_2 \vb{v}_2 \, \Rightarrow \, \vb{v}_2 = \frac{\vb{p}_2}{m_2}
\end{align*}

	allora possiamo trovare:
	
\begin{align*}
	\vb{p}_1 = \frac{m_1}{m_1 + m_2} \, \vb{P}_{\text{CM}} + \vb{p}_{\text{rel}} \qquad ; \qquad \vb{p}_2 = \frac{m_2}{m_1 + m_2} \, \vb{P}_{\text{CM}} - \vb{p}_{\text{rel}}
\end{align*}

	Se adesso si vanno a sostituire i due termini nell'Hamiltoniana di partenza, si trova che:
	
\begin{align*}
	\hat{H} = \frac{\vb*{\hat{P}}^2_{\text{CM}} }{2M} + \frac{\vb*{\hat{p}}^2_{\text{rel}} }{2\mu} + V\qty(r)
\end{align*}

	Questa è l'Hamiltoniana \textit{separata} nelle coordinate del centro di massa e relative. La prima parte:
	
\begin{align*}
	\frac{\vb*{\hat{P}}^2_{\text{CM}}}{2M}
\end{align*}

	descrive il moto del centro di massa, mentre la seconda parte:
	
\begin{align*}
	\frac{\vb*{\hat{p}}^2_{\text{rel}}}{2\mu} + V\qty(r)
\end{align*}
	
	descrive il moto relativo delle due particelle nel potenziale scelto.
	\\
	\\
	Che conseguenze ha tutto ciò?
	\\
	\\
	L'aver separato la coordinata relativa da quella del centro di massa permette di fattorizzare la funzione d'onda spaziale come:
	
\begin{align*}
	\Phi\qty(\vb{R}, \vb{r}) = \phi_{\text{CM}}\qty(R) \, \phi_{\text{rel}}\qty(r)
\end{align*}
	
	dove la funzione d'onda spaziale relativa è formalmente identica alla funzione d'onda dell'atomo di Idrogeno, perciò è possibile scrivere:
	
\begin{align*}
	\phi_{\text{rel}}\qty(r) = R_{nl}\qty(r) \, Y_l^m\qty(\theta, \phi)
\end{align*}
	
	Consideriamo adesso il caso di \textbf{particelle identiche}: cosa succede se applico la simmetria di scambio alla funzione d'onda?
	\\
	\\
	La parte dipendente dal centro di massa non subirà cambiamenti, ma la parte dipendente dalla coordinata relativa si. Infatti per scambio da $1$ a $2$ e viceversa, ho che: $r = \vb{r}_1 - \vb{r}_2 \to \vb{r}_2 - \vb{r}_1 = -r$, che non apporta cambiamenti alla funzione, e $\qty(\theta, \phi) \to \qty(\pi - \theta, \phi + \pi)$ che invece secondo la parità delle armoniche sferiche sotto inversione spaziale:
	
\begin{align*}
	\Pi Y_l^m \qty(\theta, \phi) = Y_l^m \qty(\pi - \theta, \phi + \pi) = \qty(-1)^l \, Y_l^m\qty(\theta, \phi)
\end{align*}
	
	che si porta dietro un fattore $\qty(-1)^l$, che determina in conclusione la \textbf{simmetria della funzione d'onda spaziale}.
	\\
	\\
	Si consideri un elettrone (\textit{fermione}) e si indichi con $\hat{H}^{\qty(1)}$ il suo spazio di singola particella. Lo spazio $\hat{H}^{\qty(1)}$ è un prodotto tensore tra il suo spazio orbitale ed il suo spazio di spin, ovvero:
	
\begin{align*}
	\hat{H}^{\qty(1)} = \mathcal{H}_{\text{orb}} \otimes \mathcal{H}_s
\end{align*}

	Sia $\qty{\ket{n}, \, n = 0, 1, 2, \ldots }$ una base discreta dell'autospazio $\mathcal{H}_{\text{orb}}$. Si abbiamo ora $N = 2$ elettroni, di cui consideriamo solo lo spin. 
	\\
	\\
	Quanti sono gli stati del sistema di due fermioni $\mathcal{H}^{\qty(1)} \otimes \mathcal{H}^{\qty(1)}$?
	\\
	\\
	Per rispondere a questa domanda dobbiamo conoscere la dimensione dello spazio orbitale: in genere essa è infinita, poiché i numeri quantici possono assumere infiniti valori, tuttavia scelgo di studiare cosa succede a $2$ elettroni su $2$ livelli energetici. Ovvero scelgo: $\dim \mathcal{H}_{\text{orb}} = 2$. Che dimensione ha lo spazio fermionico associato a $\mathcal{H}_-$? Poiché:
	
\begin{align*}
	\dim \mathcal{H}^{\qty(1)} = \dim \mathcal{H}_{\text{orb}} \cdot \dim \mathcal{H}_{s} = 2 \, \qty(2 s + 1) \, \equiv \, k
\end{align*}

	si ha, per $N = 2$ (prendiamo infatti in esame $2$ elettroni) che vale :
	
\begin{align*}
	\dim \mathcal{H}_- = \mqty(k \\ N) = \mqty(2\qty(2s + 1) \\ 2) 
\end{align*}

	ricordando che:
	
\begin{align*}
	\mqty(k \\ N) = \frac{k!}{N!\, \qty(k - N)!} 
\end{align*}

	allora:

\begin{align*}
	 \mqty(2\qty(2s + 1) \\ 2)  = \qty(2s + 1)\, \qty(4s + 1)
\end{align*}

	Adesso \textit{costruisco due operatori che commutano con $\hat{P}$, l'\textcolor{cyan}{energia totale del sistema} e lo \textcolor{orange}{spin totale del sistema}:}
	
\begin{align*}
	& \textcolor{cyan}{\hat{H} } = \hat{H}^{\qty(1)} \otimes \mathbb{I} + \mathbb{I} \otimes \hat{H}^{\qty(1)} \\
	& \textcolor{orange}{ \hat{S}_z }= \hat{S}_z^{\qty(1)} \otimes \mathbb{I} + \mathbb{I} \otimes \hat{S}_z^{\qty(1)}
\end{align*}

	tali che, detta la base di $\mathcal{H}_N = \mathcal{H}^{\qty(1)} \otimes \mathcal{H}^{\qty(1)}$ la seguente:
	
\begin{align*}
	\qty{ \ket{n_1}\ket{\sigma_1}, \, \ket{n_2} \ket{\sigma_2} ; \, n_i = 1, 2 \, ;  \, \sigma = -s, \ldots, + s  }
\end{align*}

	allora si ha:
	
\begin{align*}
	& \hat{H}\ket{n_1, \sigma_1; \, n_2, \sigma_2} = \qty(E_{n_1} + E_{n_2}) \ket{n_1, \sigma_1; \, n_2, \sigma_2} \\
	& \hat{S}_z\ket{n_1, \sigma_1; \, n_2, \sigma_2} = \qty(\sigma_{n_1} + \sigma_{n_2}) \ket{n_1, \sigma_1; \, n_2, \sigma_2}
\end{align*}

	quindi uno stato del sistema composto è assegnato se ho $n_1, n_2$ e $\sigma_1, \sigma_2$. In generale ho $3$ livelli energetici distinti: per un generico valore di $s$, come si distribuiscono quindi i $\qty(2s + 1) \qty(4s + 1)$ stati sui vali livelli?


\begin{itemize}
	\item $E = 2 E_1$
		
		Entrambi gli elettroni sono sul livello $E_1$. Ho quindi $n_1 = n_2 = 1$, quindi una base è $\qty{\ket{1, \sigma_1}, \ket{1, \sigma_2}}$. Questo significa che entrambi gli elettroni si trovano sul livello energetico fondamentale $E_1$ (quindi $n_1=1$ e $n_2=1$). La parte spaziale della funzione d'onda è quindi il prodotto dello stato $1$ con se stesso:
		
	\begin{align*}
		\psi_{\text{spazio}} = \ket{1}\ket{1}
	\end{align*}
	
		se scambi le due particelle ottieni ancora $\ket{1}\ket{1}$, quindi la parte spaziale è \textit{simmetrica}. I fermioni devono avere una funzione d'onda totale \textit{antisimmetrica} (devono cambiare segno se li scambi).
La funzione totale è il prodotto di Spazio $\times$ Spin. Ossia:

	\begin{align*}
		\Psi_{\text{tot}} = \psi_{\text{spazio}} \otimes \chi_{\text{spin}}
	\end{align*}

		Affinché il totale sia antisimmetrico, poiché sistema di fermioni, dato che lo spazio è simmetrico, la parte di spin deve essere obbligatoriamente \textit{antisimmetrica}. La funzione d'onda relativa a questa configurazione energetica è allora:
	
	\begin{align*}
		\Psi_{\sigma_1, \sigma_2}^{\qty(2E_1)} = \frac{1}{\sqrt{2}} \, \qty(\ket{1, \sigma_1} \ket{1, \sigma_2} - \ket{1, \sigma_2} \ket{1, \sigma_1})
	\end{align*}
	
		quanti sono questi stati?
		\\
		\\
		Ora il problema diventa puramente combinatorio: quanti stati di spin antisimmetrici esistono per due particelle con spin $s$?
		\\
		\\
		Il numero di spin per una singola particella: $N_s = 2s + 1$ e dobbiamo scegliere $2$ stati di spin diversi ($\sigma_1 \neq \sigma_2$) tra gli $N_s$ disponibili per formare una combinazione antisimmetrica del tipo $(\sigma_1 \sigma_2 - \sigma_2 \sigma_1)$. Non possono essere uguali altrimenti la differenza fa zero (Principio di Pauli), allora vale la formula?
		
	\begin{align*}
		\text{numero di conteggi} & = \mqty(2s + 1 \\ 2) = \frac{\qty(2s + 1)!}{2! \, \qty(\qty(2s + 1) - 2)!} = \frac{\qty(2s + 1) \qty(2s) \qty(2s - 1)!}{ 2! \qty( 2s - 1)!} \\
		& = \frac{\qty(2s + 1) \qty(2s)}{2} = s  \qty(2 s + 1)
	\end{align*}
	
		Ho quindi $s \qty(2s + 1)$ stati possibili su questo livello energetico.
	
	\begin{myapprf}
		\textbf{Esempio:}
		\\
		\\
		Per 2 fermioni identici, per esempio $2$ elettroni, otteniamo: $s = 1/2$. Allora abbiamo:
		\begin{align*}
			s\qty(2s + 1) = 1
		\end{align*}
		È corretto? Sì. Se due elettroni sono nello stesso orbitale spaziale (simmetrico), devono essere nel singoletto di spin (antisimmetrico):
		\begin{align*}
			\frac{1}{\sqrt{2}} \, \qty(\ket{+} \ket{-} - \ket{-} \ket{+})
		\end{align*}
		C'è un solo stato possibile. Se provassero a stare in uno stato di spin simmetrico (i 3 stati del tripletto), la funzione totale diventerebbe simmetrica, il che è vietato per i fermioni (Pauli).

	\end{myapprf}
		
	\item $E = E_1 + E_2$
	
		Un elettrone è sul livello $E_1$ mentre l'altro elettrone è sul livello $E_2$. Quindi in questo caso abbiamo un elettrone sul livello energetico fondamentale $E_1$ (chiamiamolo stato orbitale 1) e l'altro sul livello eccitato $E_2$ (stato orbitale 2).
Poiché le energie sono diverse, le due particelle occupano zone spaziali diverse (o meglio, orbitali diversi).
	\\
	\\
		\textit{Il Principio di Esclusione di Pauli dice che due fermioni non possono avere tutti i numeri quantici uguali.} Nel caso precedente ($E=2E_1$), entrambi erano nello spaziale $1$, quindi dovevano avere spin diversi per non essere identici. In questo caso, invece, i fermioni sono già distinti dalla loro posizione (uno è in $1$, l'altro è in $2$).
	\begin{itemize}
		\item L'elettrone sul livello $1$ ha numeri quantici $(n=1, \sigma_1)$.
		\item L'elettrone sul livello $2$ ha numeri quantici $(n=2, \sigma_2)$.
	\end{itemize}
	
		Dato che $n=1$ è diverso da $n=2$, i due stati completi sono sempre diversi, indipendentemente dallo spin!
Anche se gli spin fossero identici ($\sigma_1 = \sigma_2$), le particelle non sarebbero nello stesso stato perché una è al "piano terra" e l'altra al "primo piano". Poiché non c'è nessun vincolo che lega lo spin del primo a quello del secondo, abbiamo la massima libertà di scelta:
	\begin{itemize}
		\item Scelta per l'elettrone su $E_1$: Può avere qualsiasi orientazione di spin. Quante opzioni ci sono? $2s+1$.
		\item Scelta per l'elettrone su $E_2$: Può avere, indipendentemente dall'altro, qualsiasi spin. Quante opzioni ci sono? $2s+1$.
	\end{itemize}
	
		Il numero totale di combinazioni è semplicemente il prodotto delle possibilità:
	\begin{align*}
		\text{numero di conteggi} = \qty(2s + 1)\qty(2s + 1) = \qty(2s + 1)^2
	\end{align*}
	
		La funzione d'onda infine è:
	\begin{align*}
		\Psi_{\sigma_1, \sigma_2}^{\qty(E_1 + E_2)} = \frac{1}{\sqrt{2}} \, \qty(\ket{1, \sigma_1} \ket{2, \sigma_2} - \ket{2, \sigma_2} \ket{1, \sigma_1})
	\end{align*}
	
		Questa formula costruisce automaticamente uno stato antisimmetrico valido per qualsiasi coppia di valori $\sigma_1$ e $\sigma_2$, proprio perché i ket $\ket{1}$ e $\ket{2}$ sono diversi.

	\item $E = 2 E_2$
	\\
	\\
		Entrambi gli elettroni sono sul livello $E_2$. In maniera del tutto analoga al primo caso allora si ottiene che la funzione totale è:
		
	\begin{align*}
		\Psi_{\sigma_1, \sigma_2}^{\qty(2E_2)} = \frac{1}{\sqrt{2}} \, \qty(\ket{2, \sigma_1} \ket{2, \sigma_2} - \ket{2, \sigma_2} \ket{2, \sigma_1})
	\end{align*}
	
		ed ho come il primo caso il numero di stati pari a:
		
	\begin{align*}
		\text{numero di conteggi} = s \qty(2s +1 )
	\end{align*}
	
\end{itemize}

	La somma delle tre degenerazioni è:
	
\begin{align*}
	2s \qty(2s + 1) + \qty(2s + 1)^2 = \qty(2s + 1) \qty(2s + 2s +1 ) = \qty(2s + 1) \qty(4s +1) = \text{dim} \, \mathcal{H}_-
\end{align*}

	Vediamo ora molto brevemente le due statistiche citate prima. Esse costituiscono un "aggiornamento" quantistico della statistica classica di Maxwell-Boltzmann.
	
\subsection{Statistica di Bose-Einstein}

	In questo caso in un qualsiasi stato può stare un numero arbitrario di particelle. Il numero medio di particelle che occupano lo stato $k$ è dato dalla statistica di Bose-Einstein e vale:
	
\begin{align*}
	\mean{n_k} = \frac{1}{e^{\qty(\varepsilon_k - \mu)/k_B T} - 1}
\end{align*}
	
	con energia media che vale:
	
\begin{align*}
	\mean{E_k} = \mean{n_k} \varepsilon_k = \frac{\varepsilon_k}{e^{\qty(\varepsilon_k - \mu)/k_B T} - 1}
\end{align*}

	diove $\varepsilon_k$ è il $k-$esimo stato e $\mu$ è il potenziale chimico.

\subsection{Statistica di Fermi-Dirac}

	In questo la forma della funzione d'onda è:
	
\begin{align*}
	\Psi\qty(\vb{r}_1, \dots, \vb{r}_N, \sigma_1, \dots, \sigma_N) = \frac{1}{\sqrt{N!}} \sum_P \qty(-1)^p P\qty{\psi_1\qty(\vb{r}_1, \sigma_1), \dots, \psi_N\qty(\vb{r}_N, \sigma_N)}
\end{align*}

	dove la somma è estesa a tutte le permutazioni $P$ delle variabili e $p$ è la parità della permutazione. Il membro di destra dell’equazione equivale allo sviluppo del \textbf{determinante di Slater}. In cui $N$ è il numero dei fermioni nel sistema. Da ciò segue subito che se due funzioni $\psi_i$ e $\psi_j$ sono uguali fra loro, ci saranno due righe uguali nel determinante che quindi si annullerà. In un sistema di fermioni identici, come già detto, non posso assegnare la stessa $\psi$ a due fermioni. Il formalismo del determinante di Slater è valido per un sistema di fermioni non interagenti in un qualsiasi tipo di potenziale e per ogni livello energetico, a patto di fare attenzione alla sua impostazione. Il numero medio delle particelle che occupano lo stato $k$ è la statistica di Fermi-Dirac.

\begin{align*}
	\mean{n_k} = \frac{1}{e^{\qty(\varepsilon_k - \mu)/k_B T} + 1}
\end{align*}
	
	con energia media che vale:
	
\begin{align*}
	\mean{E_k} = \mean{n_k} \varepsilon_k = \frac{\varepsilon_k}{e^{\qty(\varepsilon_k - \mu)/k_B T} + 1}
\end{align*}

\section{Particelle Identiche nel Potenziale Coulombiano}

	In questa sezione vediamo degli esempi di sistemi di particelle identiche nel potenziale Coulombiano, studiandone a grandi linee il comportamento.
	
\subsection{L'Atomo di Elio non Interagente}

	Consideriamo un sistema composto da $2$ elettroni (si tratta di $2$ fermioni identici) che interagiscono con un nucleo $Z = 2$. L'Hamiltoniana del sistema è data da:
	
\begin{align*}
	\hat{H} =  \textcolor{blue}{ \frac{\vb*{\hat{p}}^2_1}{2m} - \frac{2e^2}{4\pi \varepsilon_0 r_1} } + \textcolor{ForestGreen}{ \frac{ \vb*{\hat{p}_2^2}}{2m} - \frac{2e^2}{4\pi \varepsilon_0 r_2} } + \textcolor{orange}{  \frac{e^2}{4\pi\varepsilon_0} \, \frac{1}{| \vb{r}_1 - \vb{r}_2 |} }  \, \equiv \, \textcolor{blue}{ \hat{H}_1 } +  \textcolor{ForestGreen}{ \hat{H}_2 } + \textcolor{orange}{\hat{H}_{12}}
\end{align*}	

	per il momento trascuriamo l'interazione fra i due elettroni e poniamo quindi:
	
\begin{align*}
	\hat{H} = \hat{H}_1 + \hat{H}_2
\end{align*}

	Voglio descriverne lo spettro tenendo conto del \textit{Principio di Pauli:} per fare ciò scrivo il sistema di singola particella tenendo conto anche dello spin.
	
\begin{itemize}
	\item \textbf{Livello Fondamentale} $\qty(n_1 = n_2 = 1)$
	
		Il livello fondamentale dell'elio è composto da due elettroni in $n = 1$. L'energia del livello, in unità di Bohr, è quindi:
	\begin{align*}
		E^{\qty(1, 1)} = E_1^{\qty(1)} + E_2^{\qty(1)}
	\end{align*}

		Il calcolo deriva dall'applicazione della formula sopra all'atomo di Elio, considerando l'approssimazione "non interagente".
	\\
	\\
		L'atomo è di Elio e quindi $Z = 2$ (due protoni nel nucleo). Entrambi gli elettroni sono nel livello più basso possibile, quindi $n_1 = 1$ e $n_2 = 1$. Allora possiamo usare la formula:
	\begin{align*}
		& E_1^{\qty(1)} = - \frac{Z^2}{n_1^2} = - 4 \\
		& E_2^{\qty(1)} = - \frac{Z^2}{n_2^2} = - 4
	\end{align*}

		e quindi:
	
	\begin{align*}
		E^{\qty(1, 1)} = E_1^{\qty(1)} +  E_2^{\qty(1)} = - 8
	\end{align*}
	
		dire che l'energia è $-8$ significa dire che l'energia totale dello stato fondamentale dell'Elio (senza repulsione tra elettroni) è pari a 8 volte l'energia di legame dell'atomo di Idrogeno ($8 \cdot 13.6 \, \text{eV} \approx 108.8 \, \text{eV}$).
		\\
		\\
		Per il principio di Pauli, poiché sto considerando due fermioni, lo stato del livello è descritto dalla funzione d'onda:
	\begin{align*}
		\Psi^{\qty(1, 1)} = \frac{1}{\sqrt{2}} \qty(\ket{1, 0, 0, +}_1 \ket{1, 0, 0, 0, -}_2 - \ket{1, 0, 0, -}_1 \ket{1, 0, 0, +}_2)
	\end{align*}
	
		ed è l'unico modo possibile. Da un semplice calcolo si trova allora che il momento angolare orbitale dello stato fondamentale è: $L = 0$. Essa infatti è diretta conseguenza della regola della \textit{somma dei momenti angolari}. Ecco i passaggi logici:
		
		\begin{enumerate}
			\item Siccome siamo nello stato fondamentale, vale che entrambi gli elettroni si trovano sul livello energetico più basso possibile in cui: $n_1 = 1$ e $n_2 = 1$. Il momento angolare orbitale $l$ può allora assumere solo valori interi da $0$, fino ad $n-1$. Poiché $n = 1$ allora l'unica scelta possibile per $l$ è $0$. Quindi ne segue che il primo elettrone ha momento angolare orbitale pari a $l_1 = 0$ mentre il secondo ha momento angolare orbitale pari a $l_2 = 0$.
		
			\item Per trovare $L$ dobbiamo fare la somma vettoriale dei momenti angolari, ossia:
			\begin{align*}
				\vb{L} = \vb{l}_1 + \vb{l}_2
			\end{align*}
			
				$L$ quindi può variare tra la differenza e la somma dei singoli moduli:
			\begin{align*}
				|l_1 - l_2| \leq L \leq l_1 + l_2
			\end{align*}
				
				separati di $1$, siccome $l_1 = l_2 = 0$ in questo caso, si ottiene che l'unica soluzione possibile è proprio:
			\begin{align*}
				L = 0
			\end{align*}

		\end{enumerate}

			Quindi $L = 0$ abbiamo detto, allora per garantire l'antisimmetria complessiva della funzione d'onda, gli elettroni saranno in singoletto di spin: con $S = 0$. Infatti:
		
		\begin{align*}
			\vb{\frac{1}{2}} \otimes \vb{\frac{1}{2}} = \vb{1} \oplus \vb {0}
		\end{align*}
		
			dove infatti il prodotto tensore indica che sto combinando due sistemi indipendenti con spin $1/2$. La somma diretta risultante indica che il risultato si scompone (si disaccoppia) in due sottospazi distinti e indipendenti. Tuttavia, in fisica dobbiamo scartare $S=1$ a causa del \textit{Principio di Pauli.} Questo perché siccome la parte spaziale è simmetrica per scambio di elettrone, la parte di spin deve essere antisimmetrica, la decomposizione ci dice che, per particelle identiche, il primo spazio di spin è simmetrico e poi seguono gli altri alternandosi, allora necessariamente scelgo solo $S=0$. Non perché $S=1$ non esista matematicamente, ma perché è proibito dalla configurazione spaziale in cui mi trovo. Se gli elettroni fossero in orbitali diversi (spazio antisimmetrico), allora avresti dovuto scegliere $S=1$. Ma stando "nella stessa stanza" (livello fondamentale), devono per forza avere spin opposti in configurazione di singoletto per non violare il principio di esclusione. Quindi, siccome sia $L$ che $S$ sono zero, allora lo stato fondamentale ha $\qty(L, S, J) = \qty(0, 0, 0)$, ricordando infatti che $J = L + S$.

	\item \textbf{Primo livello eccitato} $\qty(n_1 = 1, n_2 = 2)$
	
		Passiamo ora al primo livello eccitato, in cui un elettrone è in $n = 1$ mentre l'altro è in $n = 2$. L'energia in questo caso è:
		
	\begin{align*}
		E^{\qty(1, 2)} = E_1^{\qty(1)} + E_2^{\qty(2)}
	\end{align*}
	
		dove:
		
	\begin{align*}
		& E_1^{\qty(1)} = - \frac{Z^2}{n_1^2} = -\frac{4}{1} = - 4 \\
		& E_2^{\qty(2)} = - \frac{Z^2}{n_2^2} = - \frac{4}{4} = -1
	\end{align*}
	
		allora:
		
	\begin{align*}
		E^{\qty(1, 2)} = - 5
	\end{align*}

		la funzione d'onda che descrive il livello è invece:
		
	\begin{align*}
		\Psi^{\qty(1, 2)} = \frac{1}{\sqrt{2}} \, \qty(\ket{1, 0, 0, \sigma}_1 \ket{2, l', m', \sigma'}_2 - \ket{2, l', m', \sigma'}_1 \ket{1, 0, 0, \sigma}_2)
	\end{align*}

		con quindi: $ n = 1$, allora: $l = m = 0$, con $\sigma = - s, -s +1, \dots, + s$ e siccome $s = 1/2$, allora: $\sigma = -1/2, 1/2$, infatti per $s  =1/2 $, ho $2s +1$ stati, ossia in questo caso, $2$. Invece, per $ n' = 2$, $ l' = 0, 1$, $m' = 0$ per $l' = 0$ mentre $m' = -1, 0, 1$ per $l' = 1$, invece $\sigma' = -1/2, 1/2$ come prima.
		\\
		\\
		Questa funzione d'onda deve essere totalmente antisimmetrica per scambio da $1$ a $2$ e viceversa, ossia per $l', m', \sigma, \sigma'$, fissati la funzione deve essere antisimmetrica. 
		\\
		\\
		Qual è la degenerazione di questo livello? L'elettrone nello stato fondamentale ha degenerazione $2$ poiché ha solo una doppia degenerazione di spin, invece l'elettrone nel primo livello ha degenerazione $8$ poiché ha due possibili stati di spin e due possibili valori di $l'$ e quindi ha anche $4$ possibili valori di $m'$ ed in totale $8$, perciò la degenerazione totale del livello è $16$. In generale si può dire subito la degenerazione del livello tenendo conto che la degenerazione di un singolo livello è pari a $2n^2$. Più nello specifico, questi $16$ stati che $L$ e $S$ avranno?
		\\
		\\
		Si ha che:
	
	\begin{align*}
		l = 0, \,\,\, l' = 0, 1 \,\, \Rightarrow \,\, L = l \otimes l' = 0 \otimes \qty( 0 \oplus 1) = 0 \oplus 1 
	\end{align*}
		
		e vale che:

	\begin{align*}
		s = \frac{1}{2}, \,\,\, s' = \frac{1}{2} \,\, \Rightarrow \,\, S = s \otimes s' = \frac{1}{2} \otimes \frac{1}{2} = 0 \oplus 1 
	\end{align*}

		allora in questo caso, essendo su due livelli distinti gli elettroni non hanno restrizioni e non dobbiamo cancellare nessun valore. Le terne di $\qty(L, S, J)$ possibili sono:
		
	\begin{align*}
		\qty(L, S, J) & = \qty(0, 0, 0) \, \Rightarrow \, \text{deg} = 1 \\
		& = \qty(0, 1, 1) \, \Rightarrow \, \text{deg} = 3 \\
		& = \qty(1, 0, 1) \, \Rightarrow \, \text{deg} = 3 \\
		& = \qty(1, 1, 0) \, \Rightarrow \, \text{deg} = 1 \\
		& = \qty(1, 1, 1) \, \Rightarrow \, \text{deg} = 3 \\
		& = \qty(1, 1, 2) \, \Rightarrow \, \text{deg} = 5 \\
	\end{align*}
	
		le cui singole degenerazioni sono calcolate a partire da $2j + 1 $ applicato a $L$ e $S$, moltiplicate fra loro e sommate in totale, oppure al valore di $J$ e sommate in totale: si ritrova una \textit{degenerazione totale} di $16$. Lo spazio totale in cui vive $J$ è generato dal prodotto dello spazio di $L$ e dello spazio di $S$ e vale:
		
	\begin{align*}
		\mathcal{D}_L \otimes \mathcal{D}_S = \bigoplus_{J = |L - S|}^{L + S} \mathcal{D}_J
	\end{align*}

\end{itemize}

\subsection{L'atomo di Litio non Interagente}
	
	Consideriamo un sistema di $3$ elettroni ($3$ fermioni identici) che interagiscono con un nucleo con $Z = 3$. L'Hamiltoniana del sistema è data in maniera analoga da:
	
\begin{align*}
	\hat{H} = \hat{H}_1 + \hat{H}_2 + \hat{H}_3 + \hat{H}_{12} + \hat{H}_{13} + \hat{H}_{23}
\end{align*}

	per il momento trascuriamo le interazioni reciproche fra gli elettroni e poniamo quindi:

\begin{align*}
	\hat{H} = \hat{H}_1 + \hat{H}_2 + \hat{H}_3
\end{align*}

	Voglio descrivere lo spettro tenendo conto del principio di Pauli: per fare questo scrivo i sistemi di singola particella tenendo conte anche dello spin.

\begin{itemize}
	\item \textbf{Livello fondamentale} ($n_1 = n_2 = 1$ e $n_3 = 2$)
	
		Al livello fondamentale l'atomo di Litio è costituito da due elettroni in $n = 1$ ed un elettrone in $n = 2$, in modo da rispettare il principio di Pauli. L'energia del livello quindi è:
		
	\begin{align*}
		E^{\qty(1, 2, 3)} = E_1^{\qty(1)} + E_2^{\qty(1)} + E_3^{\qty(2)}
	\end{align*}
	
		e dato che non sono interagenti vale che:
		
	\begin{align*}
		& E_1^{\qty(1)} = - \frac{Z^2}{n_1^2 } = - \frac{9}{1} = - 9 \\
		& E_2^{\qty(1)} = - \frac{Z^2}{n_2^2 } = - \frac{9}{1} = - 9 \\
		& E_3^{\qty(2)} = - \frac{Z^2}{n_3^2 } = - \frac{9}{4} 
	\end{align*}
	
		e quindi:
	
	\begin{align*}
		E^{\qty(1, 1, 2 )} = E_1^{\qty(1)} + E_2^{\qty(1)} + E_3^{\qty(2)} = - 9 - 9 - \frac{9}{4} = - \frac{81}{4} = -20.25
	\end{align*}
	
		ciò significa che l’energia totale dello stato fondamentale del Litio (senza repulsione tra elettroni) è pari a $81/4$ volte l’energia di legame dell’atomo di Idrogeno.
		\\
		\\
		Che degenerazione ha questo livello? 
		\\
		\\
		A livello pratico possiamo calcolare le dimensioni dei sottospazi di singola particella antisimmetrizzati e poi moltiplicarle tra loro. In generale, \textit{al livello $n$ avrò che la dimensione del sottospazio antisimmettrizzato di $N$ fermioni è dato da:}
	
	\begin{nb}
	\begin{align*}
		\text{dim} \, \mathcal{H}_- = \mqty(n^2\qty(2s + 1) \\ N )
	\end{align*}
	\end{nb}

		quindi nel nostro caso abbiamo che:
	
	\begin{align*}
		& \text{dim} \, \mathcal{H}_-^{\qty(1, 2)} = \mqty(2 \\ 2) = 1 \\
		& \text{dim} \, \mathcal{H}_-^{\qty(3)} = \mqty(8 \\ 1) = \frac{8!}{1! \, 7!} = 8
	\end{align*}
	
		allora:
	
	\begin{align*}
		\text{dim} \, \mathcal{H}_- = 1 \cdot 8 = 8
	\end{align*}

		Per ritrovare gli stessi risultati a livello concettuale, sappiamo che a $n = 1$ ho sicuramente $S = L = 0$, l'unica $ J $ possibile è allora $0$ e quindi la degenerazione totale del livello $n = 1$ è $1$. Ad $n = 2$ invece ho solo un elettrone, allora $S = s = 1/2$ (degenerazione $2$), inoltre, essendo che siamo a $n = 2$, allora $l = 0, 1$ in questo livello ed essendo da solo l'elettrone vale che: $L = 0, 1$ (degenerazione $1$ e $3$), quindi la degenerazione totale del livello $n = 2$, vale: $\text{deg} = 2 \cdot 1 + 2 \cdot 3  = 8$. Pertanto ritroviamo che la degenerazione totale del livello fondamentale di Litio è: $1 \cdot 8 = 8$.
		\\
		\\
		Chi sono questi stati?
		\\
		\\
		Ragionando per livelli ho che:
	\begin{align*}
		& \qty(L, S)_{n = 1} = \qty(0, 0) \\
		& \qty(L, S)_{n = 2} = \qty(0, \frac{1}{2}) \oplus \qty(1, \frac{1}{2}) 
	\end{align*}
	
		allora:
		
	\begin{align*}
		\qty(L, S) = \qty(0, \frac{1}{2}) \oplus \qty(1, \frac{1}{2}) 
	\end{align*}
	
		pertanto le terne$\qty(L, S, J)$ possibili sono:
		
	\begin{align*}
		\qty(L, S, J) & = \qty(0, \frac{1}{2}, \frac{1}{2}) \, \Rightarrow \, \text{deg} = 2 \\
		& = \qty(1, \frac{1}{2}, \frac{1}{2}) \, \Rightarrow \, \text{deg} = 2 \\
		& = \qty(1, \frac{1}{2}, \frac{3}{2}) \, \Rightarrow \, \text{deg} = 4 
	\end{align*}
	
		la cui somma è $8$ come riprova del fatto.
		\\
		\\
		Qual è una base di autostati dello stato fondamentale? in termini della base singola di particella: $\ket{n, l, m, \sigma}$ ? Facciamo un primo tentativo grossolano:
		
	\begin{align*}
		\psi = \frac{1}{\sqrt{2}} \, \qty(\ket{1, 0, 0, +}_1 \ket{1, 0, 0, -}_2 - \ket{1, 0, 0, -}_1 \ket{1, 0, 0, +}_2 ) \otimes \ket{2, l, m, \pm}_3
	\end{align*}
	
		in questo modo stiamo prendendo il "nocciolo" dell'atomo (i due elettroni $1s$ che formano un singoletto, come nell'Elio) e stiamo "incollando" accanto il terzo elettrone nel livello 2. Perché questa configurazione è grossolana? Possiamo dire che questa configurazione sia non fisica perché se scambiamo l'elettrone $1$ con il $3$, la formula non ha una simmetria definita. In natura tutti gli elettroni devono essere scambiabili tra loro, questo è il significato di \textit{fermioni identici}. Quindi non possiamo semplicemente aggiungere il terzo elettrone al risultato dell'Elio. Come si fa allora? Dobbiamo prendere tutti e tre gli elettroni, metterli in un mixer matematico (il determinante di Slater) e mescolarli finché non è impossibile dire chi è chi. Allora:
		
	\begin{align*}
		\Psi_0 = \frac{1}{\sqrt{3!}} \mqty| A(1) & B(1) & C(1) \\  A(2) & B(2) & C(2) \\ A(3) & B(3) & C(3)   |
	\end{align*}
		
		dove nel nostro caso:
		
	\begin{itemize}
		\item $A = \ket{1, 0, 0, +}$ (stato $1s$ spin su)
		\item $B = \ket{1, 0, 0, -}$ (stato $1s$ spin giù)
		\item $C = \ket{2, l, m, \pm}$ (stato eccitato $n  =2$)
	\end{itemize}
	
		la matrice si legge nel seguente modo: le \textit{righe} rappresentano le particelle (elettrone $1, 2, 3$) mentre le \textit{colonne} rappresentano gli stati disponibili. Se provassimo a scambiare due particelle (es. riga 1 con riga 2), per le proprietà dei determinanti, il risultato cambia segno. Questo soddisfa perfettamente il Principio di Pauli per tutti e tre gli elettroni contemporaneamente. Allora possiamo scrivere direttamente:
		
	\begin{align*}
		\Psi_0 = \frac{1}{\sqrt{3!}} \mqty| \ket{1, 0, 0, +}_1 & \ket{1, 0, 0, -}_1 & \ket{2, l, m, \pm}_1 \\ \ket{1, 0, 0, +}_2 & \ket{1, 0, 0, -}_2 & \ket{2, l, m, \pm}_2 \\  \ket{1, 0, 0, +}_3 & \ket{1, 0, 0, -}_3 & \ket{2, l, m, \pm}_3   |
	\end{align*}

		Il determinante di Slater può anche essere scritto come:
		
	\begin{align*}
		\Psi_0 = \frac{1}{\sqrt{3!}} \sum_{ijk} \varepsilon_{ijk} \qty(\ket{1, 0, 0, +}_i \ket{1, 0, 0, -}_j - \ket{1, 0, 0, -}_i \ket{1, 0, 0, +}_j) \otimes \ket{2, l, m, \pm}_k
	\end{align*}	
	
		dove si ricorda che $\varepsilon_{ijk}$ è un operatore che vale:
	
	\begin{itemize}
		\item $+1$, se la permutazione è pari ($1,2,3$ e cicliche),
		\item $0$, se due indici sono uguali,
		\item $-1$,  se la permutazione è dispari ($2, 1, 3$ e cicliche).
	\end{itemize}

	\item \textbf{Primo livello eccitato} ($n_1 = n_2 = 1$ e $n_3 = 3$)
		
		Il primo livello eccitato è quella in cui $2$ elettroni sono in $n = 1$ ed il terzo in $n = 3$, poiché la configurazione in cui un altro elettrone passa dal livello $n = 1$ al livello $ n = 2$ ha energia più alta.
		\\
		\\
		L'energia del livello è: 
		
	\begin{align*}
		E^{\qty(1, 1, 3)} = E_1^{\qty(1)} + E_2^{\qty(1)} + E_3^{\qty(3)} = \qty(-1-1-\frac{1}{9}) Z^2 = -19
	\end{align*}
	
		infatti il primo livello eccitato non potrebbe essere un elettrone in $n = 1$ e due in $n= 2$, poiché l'energia sarebbe più alta ossia:
		
	\begin{align*}
		E^{\qty(1, 2, 2)} = E_1^{\qty(1)} + E_2^{\qty(2)} + E_3^{\qty(2)} = \qty(-1-\frac{1}{4}-\frac{1}{4}) Z^2 = -13.5
	\end{align*}
	
		Analizziamo ora la degenerazione del livello, otteniamo che:
 
 	\begin{align*}
		\text{dim} \, \mathcal{H}_- = \mqty(2 \\ 2) \, \mqty(18 \\ 1) = 18
	\end{align*}
	
		Ragionando per livelli ho che:
		
	\begin{align*}
		& \qty(L, S)_{n = 1} = \qty(0, 0) \\
		& \qty(L, S)_{n=3} = \qty(0, \frac{1}{2}) \oplus \qty(1, \frac{1}{2}) \oplus \qty(2, \frac{1}{2}) 
	\end{align*}

		allora:
		
	\begin{align*}
		\qty(L, S) & = \qty(L, S)_{n = 1} \otimes \qty(L, S)_{n = 2} \\
		& = \qty(0, 0) \otimes \qty( \qty(0, \frac{1}{2}) \oplus \qty(1, \frac{1}{2}) \oplus \qty(2, \frac{1}{2}) ) = \qty(0, \frac{1}{2}) \oplus \qty(1, \frac{1}{2}) \oplus \qty(2, \frac{1}{2}) 
	\end{align*}

		pertanto le terne $\qty(L, S, J)$ possibili sono:
		
	\begin{align*}
		\qty(L, S, J) & = \qty(0, \frac{1}{2}, \frac{1}{2}) \, \Rightarrow \,  \text{deg} = 2 \\
		& = \qty(1, \frac{1}{2}, \frac{1}{2}) \, \Rightarrow \,  \text{deg} = 2 \\
		& = \qty(1, \frac{1}{2}, \frac{3}{2}) \, \Rightarrow \,  \text{deg} = 4 \\
		& = \qty(2, \frac{1}{2}, \frac{3}{2}) \, \Rightarrow \,  \text{deg} = 4 \\
		& = \qty(2, \frac{1}{2}, \frac{5}{2}) \, \Rightarrow \,  \text{deg} = 6
	\end{align*}
	
		la degenerazione totale è $18$ com'è giusto che sia.
			
\end{itemize}

\subsection{L'atomo di Berillio non Interagente}

	Consideriamo, infine, un sistema di $4$ elettroni (4 fermioni identici) che interagiscono con un nucleo con $Z = 4$. L'Hamiltoniana del sistema è data in maniera analoga da:
	
\begin{align*}
	\hat{H} = \hat{H}_1 + \hat{H}_2 + \hat{H}_3 + \hat{H}_4 + \hat{H}_{\text{int}}
\end{align*}
	
	per il momento trascuriamo le interazioni reciproche fra gli elettroni e poniamo: $\hat{H} = \hat{H}_1 + \hat{H}_2 + \hat{H}_3 + \hat{H}_4$. Voglio descrivere lo spettro tenendo conto del principio di Pauli. Per fare ciò scrivo i sistemi di singola particella tenendo conto dello spin.
	
\begin{itemize}
	\item \textbf{Livello fondamentale} ($n_1 = n_2 = 1$ e $n_3 = n_4 = 2$)
	
		Al livello fondamentale l'atomo di Berillio è costituito da due elettroni in $n = 1$ e da due elettroni in $n = 2$, in modo da rispettare il principio di Pauli. L'energia del livello quindi, in unità di Bohr, è:

\begin{align*}
	E^{\qty(1, 1, 2, 2)} = E_1^{\qty(1)} + E_2^{\qty(1)} + E_3^{\qty(2)} + E_4^{\qty(2)} = \qty(-1-1-1/4-1/4) Z^2 = -40
\end{align*}
	
		che degenerazione ha questo livello? 
		\\
		\\
		A livello pratico possiamo calcolare le dimensioni dei sottospazi di singola particella antisimmetrizzati e poi moltiplicarle tra loro. La degenerazione è:
		
	\begin{align*}
		\text{dim} \, \mathcal{H}_- = \mqty(2 \\ 2) \, \mqty(8 \\ 2 ) = 28
	\end{align*}
	
		Chi sono questi stati?
		\\
		\\
		Ragioniamo: i due elettroni sul livello $n = 1$, sicuramente devono rispondere al principio di Pauli, poiché l'unico valore possibile del momento angolare orbitale è $l = 0$ di conseguenza:
		
	\begin{align*}
		\qty(L, S)_{n = 1} = \qty(0, 0)
	\end{align*}

		Siccome i primi due elettroni riempiono il core, ci concentriamo sui due di valenza (come visto nel caso del Litio). Dunque guardiamo il secondo livello. per quanto riguarda gli elettroni sul livello $n = 2$ il discorso è diverso infatti: i valori possibili del momento angolare orbitale sono $l = 0, 1$ (poiché $n = 2$) e di conseguenza ho:
		
	\begin{align*}
		L = \qty(1 \oplus 0) \otimes \qty(1 \oplus 0) = \textcolor{red}{(0)} \oplus \textcolor{blue}{(1)} \oplus \textcolor{magenta}{(1)} \oplus \textcolor{orange}{\qty( 0 \oplus 1 \oplus 2)}
	\end{align*}

		dove:
		
	\begin{itemize}
		\item $\textcolor{red}{(0)} \, \Rightarrow \,$ viene da: $l_1 = 0$ e $l_2 = 0$;
		\item $\textcolor{blue}{(1)} \, \Rightarrow \,$ viene da: $l_1 = 0$ e $l_2 = 1$;
		\item $\textcolor{magenta}{(1)} \, \Rightarrow \,$ viene da: $l_1 = 1$ e $l_2 = 0$;
		\item $\textcolor{orange}{\qty( 0 \oplus 1 \oplus 2)} \, \Rightarrow \,$ viene da: $l_1 = 1$ e $l_2 = 1$.
	\end{itemize}

		quindi dobbiamo stare attenti ed applicare il \textit{principio di Pauli} solo a quei valori di $L$ che sono dati da valori uguali di $l_1$ ed $l_2$. 
		\\
		\\
		Poiché $s_1 = s_2 = 1/2$, allora abbiamo che:
		
	\begin{align*}
		S = 0 \oplus 1
	\end{align*}
	
		e quindi seguendo la parità spaziale $\qty(-1)^l$ ho che i $28$ stati sono dati dalle coppie che seguono il principio di Pauli ($\qty(1s)^2$ e $\qty(2p)^2$) e dalle coppie prese come combinazione dei valori di $L$ ed $S$ (orbitale $\qty(1s)^1 \qty(2p)^1$)).
		\\
		\\
		Quindi ho che, quando gli $l$ sono diversi, gli elettroni occupano orbitali diversi, ovvero sono elettroni non equivalenti e quindi vanno presi tutti, come possiamo notare dai valori possibili di $L$, per $l$ diversi, si ottiene sempre $L = 1$, $S$ invece sarà $S = 1/2 \otimes 1/2 = 1, 0$, allora ($J = L \otimes S$):
		
	\begin{align*}
		\qty(L, S, J) & = \qty(1, 0, 1) \, \Rightarrow \, \text{deg} = 3 \\ 
		& = \qty(1, 1, 0) \, \Rightarrow \, \text{deg} = 1 \\
		& = \qty(1, 1, 1) \, \Rightarrow \, \text{deg} = 3 \\
		& = \qty(1, 1, 2) \, \Rightarrow \, \text{deg} = 5 
	\end{align*}
	
		la degenerazione quando gli orbitali sono diversi è la somma delle degenerazioni appena trovate, quindi $12$. Controlliamo ora quelli in cui $l$ sono uguali, quindi quando: $l_1 = l_2  = 0$ e quindi un solo $L = 0$, $S = 0,1$ ma siccome sta volta gli elettroni sono equivalenti, ne segue che devo controllare la simmetria, siccome $L = 0$ è simmetrico, lo spin deve essere antisimmetrico, quindi butto la soluzione con $S = 1$, antisimmetrica, e tengo solo quella simmetrica ($S = 0$):
		
	\begin{align*}
		\qty(L, S, J) = \qty(0, 0, 0) \, \Rightarrow \, \text{deg} = 1
	\end{align*}

		per invece $l_1 = l_2 = 1$ si ottiene $L = 0, 1, 2$ e $S = 0, 1$, Pauli seleziona solo:
	
	\begin{align*}
		\qty(L, S, J) & = \qty(0, 0, 0) \, \Rightarrow \, \text{deg} = 1 \\ 
		& = \qty(2, 0, 2) \, \Rightarrow \, \text{deg} = 5 \\ 
		& = \qty(1, 1, 0) \, \Rightarrow \, \text{deg} = 1 \\
		& = \qty(1, 1, 1) \, \Rightarrow \, \text{deg} = 3 \\
		& = \qty(1, 1, 2) \, \Rightarrow \, \text{deg} = 5 \\
	\end{align*}
		
		quindi quando gli $l$ sono diversi, si ottengono $16$ stati. 
		\\
		\\
		La somma totale è quindi $28$ proprio come volevamo.
		\\
		\\
		Cerchiamo ora una base di autostati dello stato fondamentale in termini della base di singola particella $\ket{n, l, m, \sigma}$. L'idea è quella di trovare la funzione $\psi_{n = 1}$ relativa ad $n = 1$ e la funzione $\psi_{n = 2}$ relativa a $n = 2$ e metterle successivamente in un prodotto tensore come funzione d'onda del \textit{ground state} del berillio:
		
	\begin{align*}
		\Psi_{\text{G.S}} = \psi_{n = 1} \otimes \psi_{n = 2}
	\end{align*}
	
		Andando per livelli:
		
	\begin{itemize}
		\item $n = 1$
		
			In questo livello come abbiamo già dimostrato altre volte, l'unica configurazione possibile è un singoletto di spin e pertanto:
			
		\begin{align*}
			\psi_{n = 1} = \frac{1}{\sqrt{2}} \qty( \ket{1, 0, 0, +}_1 \ket{1, 0, 0, -}_2 - \ket{1, 0, 0, -}_1 \ket{1, 0, 0, +}_2)
		\end{align*}

		\item $n = 2$
		
			In questo livello abbiamo detto che tutte le configurazioni possibili sono: $\qty(L, S) = \qty(0, 0) \qty(1, 0) \qty(1, 1) \qty(2, 0)$, pertanto a seconda del fatto se mi trovo nello stato di singoletto o di tripletto la funzione d'onda spaziale deve essere rispettivamente simmetrica o antisimmetrica. Consideriamo gli stati non fisici per $n = 2$ per lo stato di singoletto $\qty(S)$ e di tripletto $\qty(T)$:
			
		\begin{align*}
			\psi_2^{S} = \frac{1}{\sqrt{4}} \qty(\ket{+}_3\ket{-}_4 - \ket{-}_3 \ket{+}_4) \otimes \qty( \ket{2, l', m'}_3 \ket{2, l'', m''}_4 + \ket{2, l'', m''}_3 \ket{2, l', m'}_4 )
		\end{align*}

			dunque:
			
		\begin{align*}
			\psi_2^{T} = \begin{cases}
    \frac{1}{\sqrt{4}} (\qty(\ket{+}_3\ket{-}_4 + \ket{-}_3 \ket{+}_4) \otimes \qty( \ket{2, l', m'}_3 \ket{2, l'', m''}_4 - \ket{2, l'', m''}_3 \ket{2, l', m'}_4 )) \\
    \frac{1}{\sqrt{6}} \ket{\pm}_3 \ket{\pm}_4 \otimes \qty( \ket{2, l', m'}_3 \ket{2, l'', m''}_4 + \ket{2, l'', m''}_3 \ket{2, l', m'}_4 ))
\end{cases}
		\end{align*}
			
	\end{itemize}

		Lo stato complessivo è:
	
	\begin{align*}
		\Psi_{\text{G.S}} = \psi_{n = 1} \otimes \psi_{n = 2}
	\end{align*}

		dovrà essere totalmente antisimmetrico per scambi di $1, 2, 3, 4$. In forma compatta la funzione d'onda relativa allora stato fondamentale dell'atomo di Berilio si scrive con il determinante di Slater nel seguente modo:	
		
	\begin{align*}
		\Psi_0 = \frac{1}{\sqrt{4!}} \mqty|  \ket{1, 0, 0, +}_1  &  \ket{1, 0, 0, -}_1 &  \ket{2, l', m', \sigma'}_1   &  \ket{2, l'', m'', \sigma''}_1   \\   \ket{1, 0, 0, +}_2  &  \ket{1, 0, 0, -}_2 &  \ket{2, l', m', \sigma'}_2   &  \ket{2, l'', m'', \sigma''}_2   \\  \ket{1, 0, 0, +}_3  &  \ket{1, 0, 0, -}_3 &  \ket{2, l', m', \sigma'}_3   &  \ket{2, l'', m'', \sigma''}_3   \\    \ket{1, 0, 0, +}_4  &  \ket{1, 0, 0, -}_4 &  \ket{2, l', m', \sigma'}_4   &  \ket{2, l'', m'', \sigma''}_4  |
	\end{align*}
	
	\item \textbf{Primo livello eccitato} ($n_1 = n_2 = 1$ e $n _3 = 2$ e $n_4 = 3$)
	
		Al primo livello eccitato l'atomo di Berilio è costituito da due elettroni in $n = 1$, un elettrone in $n = 2$ ed un altro in $n = 3$. L'energia del livello quindi, in unità di Bohr, è:
		
	\begin{align*}
		E^{\qty(1, 1, 2, 3)} = E_1^{\qty(1)} + E_2^{\qty(1)} + E_3^{\qty(2)} + E_4^{\qty(3)} = \qty(-1-1-\frac{1}{4}-\frac{1}{9}) Z^2 = - \frac{340}{9}
	\end{align*}

		Che degenerazione ha questo livello?
		\\
		\\
		A livello pratico possiamo calcolare le dimensioni dei sottospazi di singola particella antisimmetrizzati e poi moltiplicarle e poi moltiplicarle tra loro. La degenerazione è:
	
	\begin{align*}
		\text{dim} \, \mathcal{H}_- = \mqty(2 \\ 2) \, \mqty( 8 \\ 1) \, \mqty(18 \\ 1) = 144
	\end{align*}

	Chi sono questi stati? Ragioniamo: i due elettroni sul livello $n = 1$ sicuramente devono rispondere al principio di Pauli, poiché l’unico valore possibile del momento angolare orbitale è $l = 0$ e, di conseguenza, proprio come prima vale che:
	
	\begin{align*}
		\qty(L, S)_{n = 1} = \qty(0, 0)
	\end{align*}
	
		per i livelli sucessivi allora si ottiene:
		
	\begin{itemize}
		\item $n = 2$
		
		\begin{align*}
			\qty(L, S)_{n = 2} = \qty(0, \frac{1}{2}) \oplus \qty(1, \frac{1}{2})
		\end{align*}

		\item $n = 3$
			
		\begin{align*}
			\qty(L, S)_{n = 3} = \qty(0, \frac{1}{2} ) \oplus \qty(1, \frac{1}{2} ) \oplus \qty(2, \frac{1}{2})
		\end{align*}
	\end{itemize}

		Allora si ottiene che:
	
	\begin{align*}
		\qty(L, S) & =  \qty(0, 0) \otimes \qty(\qty(0, \frac{1}{2} ) \oplus \qty(1, \frac{1}{2} )) \otimes \qty( \qty(0, \frac{1}{2} ) \oplus \qty(1, \frac{1}{2} ) \oplus \qty(2, \frac{1}{2})) \\
		& = \qty(\qty(0, \frac{1}{2} ) \oplus \qty(1, \frac{1}{2} )) \otimes \qty( \qty(0, \frac{1}{2} ) \oplus \qty(1, \frac{1}{2} ) \oplus \qty(2, \frac{1}{2})) \\
		& = \qty(0, 0) \oplus \qty(0, 1) \oplus \qty(1, 0) \oplus \qty(1, 1) \dots
	\end{align*}
\end{itemize}

		si trovano tutti calcolandoli in questo modo (per esempio):
	
	\begin{align*}
		\qty(1, \frac{1}{2} ) \otimes \qty(1, \frac{1}{2}) = \qty(0, 0) \oplus \qty(0, 1) \oplus \qty(1, 0) \oplus \qty(1, 1) \oplus \qty(2, 0) \oplus \qty(2, 1)
	\end{align*}
 
 		possiamo fare tutti i calcoli, come vedi sono molti, più di $40$ in tutto e se calcoli rispetti $J$ con degenerazione il conto alla fine torna perfettamente.
		
\section{Particelle Identiche nel Potenziale Armonico}

	In questa sezione ora vediamo degli esempi di particelle identiche in un potenziale armonico tridimensionale, studiandone a grandi linee il comportamento.

\subsection{Sistema di $2$ Fermioni Identici nel Potenziale Armonico}

	Consideriamo ora due elettroni nel potenziale armonico:
	
\begin{align*}
	\hat{H}_0 = \hat{H}^{\qty(1)} + \hat{H}^{\qty(2)}
\end{align*}
	
	con:
	
\begin{align*}
	\hat{H}^{\qty(i)} = \frac{\hat{p}^2_i}{2m} + \frac{1}{2} m \omega^2 r_i^2
\end{align*}
	
	dove $\hat{H}^{\qty(i)}$ è l'Hamiltoniana di singola particella. Sia la base di singola particella: $\qty{\ket{n, l, m, \sigma}}$.
	
\begin{itemize}
	
	\item \textbf{Livello fondamentale} ($n_1 = 0, n_2 = 0$ e $l_1 = 0, l_2 = 0$)
	
		Al livello fondamentale ho due elettroni in $n = 0$, quindi l'energia del livello sarà:
		
	\begin{align*}
		E_0 = E_{0}^{\qty(1)} + E_0^{\qty(2)}  
	\end{align*}
	
		ricorda che l'Hamiltoniana di un oscillatore armonico 3D isotropo (stessa frequenza $\omega$ in tutte le direzioni) è la somma di tre oscillatori 1D indipendenti lungo gli assi $x, y, z$, allora vale che:
	
	\begin{align*}
		E = \qty(n + \frac{3}{2}) \hbar \omega
	\end{align*}
		
		allora allo stato fondamentale:

	\begin{align*}
		E_0 = E_{0}^{\qty(1)} + E_0^{\qty(2)}  = \qty(0 + \frac{3}{2}) \hbar \omega + \qty(0 + \frac{3}{2}) \hbar \omega = 3 \hbar \omega
	\end{align*}

		calcoliamo ora la degenerazione, in questo caso armonico, per trovare la dimensione dello spazio antisimmetrico $\mathcal{H}_-$, si usa la solita regola generale combinatoria per i fermioni:
		
	\begin{align*}
		\text{dim} \, \mathcal{H}_- = \mqty(\text{stati totali disponibili} \\ \text{numero di particelle}) = \mqty(k \\ N)
	\end{align*}

		come facciamo a capire $k$ quanto vale? Sappiamo infatti che $N = 2$, ma $k$?
		\\
		\\
		In generale vale che:
		
	\begin{align*}
		k = g_n \, g_s
	\end{align*}
	
		dove $g_n$ è la degenerazione spaziale, mentre $g_s$ è la degenerazione di spin. Per un oscillatore armonico $3D$ isotropo la degenerazione del livello $n$ è data da:
		
	\begin{align*}
		g_n = \frac{\qty(n+1)\qty(n+2)}{2}
	\end{align*}
	
		quindi per $n = 0$, $g_n = 1$. Mentre la degenerazione di spin, per un elettrone vale, siccome $s = 1/2$:
		
	\begin{align*}
		g_s = 2s + 1 = 2
	\end{align*}
	
		allora:
		
	\begin{align*}
		k = 2
	\end{align*}

		e quindi:
	
	\begin{align*}
		\text{dim} \, \mathcal{H}_- =\mqty(k \\ N) = \mqty(2 \\ 2) = 1
	\end{align*}
	
		Il risultato "$1$" significa che c'è un solo modo possibile di costruire lo stato fondamentale: dobbiamo mettere entrambi gli elettroni nell'unico orbitale spaziale disponibile, e per il principio di Pauli devono avere obbligatoriamente spin opposti (Singoletto). Quindi come ci aspettavamo il livello fondamentale è non degenere.
		\\
		\\
		Qual è una base ortonormale di questo livello in termini di basi di singola particella?
		\\
		\\
		Componendo i momenti angolari abbiamo che:
		
	\begin{align*}
		& L = l_1 \otimes l_2 = 0 \otimes 0 = 0 \\
		& S = s_1 \otimes s_2 = \frac{1}{2} \otimes \frac{1}{2} = 1 \oplus 0
	\end{align*}
	
		poiché ho un sistema di due fermioni identici, la funzione d'onda complessiva del livello deve essere antisimmetrica per scambio, perciò, seguendo la parità $\qty(-1)^l$ per il momento angolare orbitale (quindi $\psi_{\text{orb}}$ è simmetrica) e considerato che $S = 0$ (singoletto) è antisimmetrico e $S = 1$ (tripletto) è simmetrico, l'unica coppia antisimmetrica possibile è $(L,S) = (0,0)$, da cui segue immediatamente che:
		
	\begin{align*}
		\qty(L, S, J) = \qty(0, 0, 0) \, \Rightarrow \, \text{non degenere}
	\end{align*}
	
		e pertanto vale che
		
	\begin{align*}
		\Psi_0 = \psi_{\text{orb}} \otimes \psi_{\text{spin}} & = \ket{0, 0, 0} \otimes \frac{1}{\sqrt{2}} \qty( \ket{0, 0, 0, +}_1 \ket{0, 0, 0, -}_2 -  \ket{0, 0, 0, -}_1  \ket{0, 0, 0, +}_2) \\
		& = \frac{1}{\sqrt{2}} \qty( \ket{0, 0, 0, +}_1 \ket{0, 0, 0, -}_2 -  \ket{0, 0, 0, -}_1  \ket{0, 0, 0, +}_2) 
 	\end{align*}

	\item \textbf{Primo livello eccitato} ($n_1 = 0$, $n_2 = 1$ e $l_1 = 0$, $l_2 = 1$)
	
		Al primo livello eccitato un elettrone si trova in $n = 0$ e l'altro in $n = 1$, quindi l'energia del livello sarà:
		
	\begin{align*}
		E = \hbar \omega \qty(\frac{3}{2} + \frac{5}{2} ) = 4 \hbar \omega
	\end{align*}
	
		con degenerazione pari a:
		
	\begin{align*}
		\text{dim} \, \mathcal{H}_- = \mqty(2 \\ 1) \, \mqty(6 \\ 1) = 12
	\end{align*}

		quindi il primo livello eccitato è $12$ volte degenere.
		\\
		\\
		Qual è la base ortonormale di questo livello in termini di basi di singola particella?
		\\
		\\
		Componendo i momenti angolari ho:
		
	\begin{align*}
		& L = l_1 \otimes l_2 = 0 \otimes 1 = 1\\
		& S = s_1 \otimes s_2 = \frac{1}{2} \otimes \frac{1}{2} = 1 \oplus 0
	\end{align*}
	
		allora:
	
	\begin{align*}
		J = 1 \otimes \qty( 1 \oplus 0) = \qty( 2 \oplus 1 \oplus 0) \oplus \qty(1)
	\end{align*}
	
		In questo caso la funzione d'onda spaziale può essere sia simmetrica che antisimmetrica, perciò ho diverse espressioni della funzione d'onda complessiva a seconda del fatto se mi trovo in uno stato di singoletto o in uno stato di tripletto di spin. In generale la funzione d’onda sarà antisimmetrica per scambio:
	
	\begin{align*}
		\psi = \frac{1}{\sqrt{2}} \qty( \ket{0, 0, 0, \sigma}_1 \ket{1, 1, m', \sigma'}_2 -  \ket{1, 1, m', \sigma'}_1 \ket{0, 0, 0, \sigma}_2 )
	\end{align*}
	
		Più esplicitamente:
		
	\begin{itemize}
	
		\item $S = 0$ (stato di singoletto): in questo caso la funzione d'onda spaziale deve essere simmetrica, poiché quella di spin è antisimmetrica:
		
		\begin{align*}
			\Psi_{S = 0} = \frac{1}{\sqrt{2}} \qty(\ket{0, 0, 0}_1 \ket{1, 1, m'}_2 + \ket{1, 1, m'}_1 \ket{0, 0, 0}_2) \otimes \frac{1}{\sqrt{2}} \qty(\ket{+}_1 \ket{-}_2 - \ket{-}_1 \ket{+}_2)
		\end{align*}
		
		\item $S = 1$ (stato di tripletto): in questo caso la funzione d'onda spaziale deve essere antisimmetrica poiché quella di spin è simmetrica. Siccome ho tre stati di spin, distinguo in base al valore di $\hat{S}_z$:
		
		\begin{align*}
			\Psi_{S = 1}^{S_z = \pm 1} = \frac{1}{\sqrt{2}} \qty(\ket{0, 0, 0}_1 \ket{1, 1, m'}_2 - \ket{1, 1, m'}_1 \ket{0, 0, 0}_2) \otimes \ket{\pm}_1 \ket{\pm}_2
		\end{align*}
		
			e
		\begin{align*}
			\Psi_{S = 1}^{S_z = 0} = \frac{1}{\sqrt{2}} \qty(\ket{0, 0, 0}_1 \ket{1, 1, m'}_2 - \ket{1, 1, m'}_1 \ket{0, 0, 0}_2) \otimes  \frac{1}{\sqrt{2}} \qty(\ket{+}_1 \ket{-}_2 + \ket{-}_1 \ket{+}_2)
		\end{align*}
	\end{itemize}
\end{itemize}

\subsection{Sistema di $3$ Fermioni Identici nel Potenziale Armonico}
	
	Considero $3$ elettroni nel potenziale armonico:
	
\begin{align*}	
	\hat{H}_0 = \hat{H}^{\qty(1)} + \hat{H}^{\qty(2)} + \hat{H}^{\qty(3)}
\end{align*}

	con, come sappiamo:
	
\begin{align*}
	\hat{H}^{\qty(i)} = \frac{\hat{p}^2_i}{2m} + \frac{1}{2} m \omega^2 r^2_i
\end{align*}

	dove $\hat{H}^{\qty(i)}$ è l'Hamiltoniana di singola particella. Sia la base di singola particella $\qty{\ket{n, l, m, \sigma}}$.
	\\
	\\
	Analizziamo allora:
	
\begin{itemize}

	\item \textbf{Livello fondamentale} ($n_1 = n_2 = 0, n_3 = 1$ e $l_1 = l_2 = 0, l_3 = 1$)
	
	\begin{myapprf}
		\textit{Nel caso di potenziale armonico} se dato un certo valore di $n$, controlla se esso è pari: $\qty(0, 2, 4, \dots)$ o dispari: $\qty(1, 3, 5, \dots)$. Se $n$ è pari, $l$ assume tutti i valori pari fino ad $n$ stesso. Se $n$ è dispari invece, $l$ assume tutti i valori dispari fino ad $l$ stesso. Facciamo un esempio: se $n$ fosse uguale a $3$, allora $l = 1, 3$, se invece fosse uguale ad $n = 1$, allora l'unico $l$ possibile è $1$ stesso. Per quanto riguarda $m$ invece valgono le stesse regole dell'atomo di idrogeno e quindi $-l \leq m \leq l$
	\end{myapprf}

	
		Al livello fondamentale ho $2$ elettroni in $n = 0$ ed uno in $n = 1$, quindi l'energia del livello sarà:
		
	\begin{align*}
		E_0 = E_0^{\qty(1)} + E_0^{\qty(2)} + E_1^{\qty(3)} = 2 \hbar \omega \qty(0 + \frac{3}{2}) + \hbar \omega \qty( 1 + \frac{3}{2}) = \frac{11}{2} \hbar \omega
	\end{align*}

		con degenerazione pari a:
		
	\begin{align*}
		\text{dim} \, \mathcal{H}_- = \mqty(2 \\ 2) \, \mqty(6 \\ 1) = 6
	\end{align*}

		il livello fondamentale ha degenerazione $6$.
		\\
		\\
		Qual è una base ortonormale di questo livello in termini di basi di singola particella? Componendo i momenti angolari, distinguo per livelli. Ho allora che:
		
	\begin{itemize}
		\item $n = 0$:
		
		\begin{align*}
			& L_{12} = l_1 \otimes l_2 = 0 \otimes 0 = 0 \\
			& S_{12} = s_1 \otimes s_2 = \frac{1}{2} \otimes \frac{1}{2} = 0 \oplus 1
		\end{align*}
		
			$S_{12}$ va scartato poiché $L_{12} = 0$ genera uno spazio simmetrico e lo spazio totale lo vogliamo antisimmetrica, allora:
			
		\begin{align*}
			\qty(L, S, J)_{12} = \qty(0, 0, 0)
		\end{align*}

		\item $n = 1$
			
		\begin{align*}
			& L_3 = l_3 = 1\\
			& S_3 = s_3 = \frac{1}{2} 
		\end{align*}
			
			allora:
			
		\begin{align*}
			\qty(L, S, J)_3 = \qty(1, \frac{1}{2}, \frac{1}{2}) \oplus \qty(1, \frac{1}{2}, \frac{3}{2})
		\end{align*}
	\end{itemize}
	
		Da cui allora otteniamo che:
		
	\begin{align*}
		\qty(L, S, J) = \qty(1, \frac{1}{2}, \frac{1}{2}) \oplus \qty(1, \frac{1}{2}, \frac{3}{2})
	\end{align*}
	
		Poiché ho un sistema di $3$ fermioni identici, la funzione d'onda complessiva del livello deve essere antisimmetrica per scambio, perciò, considero lo stato (non fisico):
		
	\begin{align*}
		\frac{1}{\sqrt{2}} \, \qty(\ket{0, 0, 0, +}_1 \ket{0, 0, 0, -}_2 - \ket{0, 0, 0, -}_1 \ket{0, 0, 0, +}_2) \otimes \ket{1, 1, m, \sigma}_3
	\end{align*}

		Gli stati permessi dal principio di Pauli sono quelli antisimmetrici per scambio di $1, 2, 3$, in forma compatta:
		
	\begin{align*}
		\psi = \frac{1}{\sqrt{3!}} \sum_{ijk} \qty(\ket{0, 0, 0, +}_i \ket{0, 0, 0, -}_j - \ket{0, 0, 0, -}_i \ket{0, 0, 0, +}_j) \otimes \ket{1, 1, m, \pm}_k
	\end{align*}

	\item \textbf{Primo livello eccitato} 
	
		Al primo livello eccitato ho $3$ possibili configurazioni:
		
	\begin{enumerate}
		\item \textit{Configurazione orbitale} $\qty(1s)\qty(2p)^2$: 
		
			In questa configurazione ho un elettrone in $n = 0$ e due elettroni in $n = 1$. In questo caso allora ho che vale:
			
		\begin{align*}
			E_1 = \frac{13}{2} \hbar \omega
		\end{align*}
		
			per comporre allora i momenti angolari distinguo in base ai livelli $n$:
		
		\begin{itemize}
			\item $n= 0$
			\begin{align*}
				\qty(L, S)_1 = \qty(0, \frac{1}{2})
			\end{align*}
			
			\item $n = 1$
			\begin{align*}
				& L_{23} = l_2 \otimes l_3 = 1 \otimes 1 = 2 \oplus 1 \oplus 0 \\
				& S_{23} = s_2 \otimes s_3 = \frac{1}{2} \otimes \frac{1}{2} = 1 \oplus 0
			\end{align*}
			
				allora, per antisimmetria:
			\begin{align*}
				\qty(L, S)_{23} = \qty(0, 0) \oplus \qty(1, 1) \oplus \qty(2, 0)
			\end{align*}

		\end{itemize}

			Quindi in questa configurazione:
			
		\begin{align*}
			\qty(L, S) & = \qty(0, \frac{1}{2}) \otimes \qty(\qty(0, 0) \oplus \qty(1, 1) \oplus \qty(2, 0)) \\
			& = \qty(0, \frac{1}{2}) \oplus \qty(1, \frac{1}{2}) \oplus \qty(1, \frac{3}{2}) \oplus \qty(2, \frac{1}{2})
		\end{align*}
		
			in questo modo ho che:
			
		\begin{align*}
			J = \frac{1}{2} \oplus \frac{1}{2} \oplus \frac{3}{2} \oplus \frac{1}{2} \oplus \frac{3}{2} \oplus \frac{5}{2} \oplus \frac{3}{2} \oplus \frac{5}{2} \, \Rightarrow \, \text{Avrò degenerazione totale pari a:} \, 30
		\end{align*}

		\item \textit{Configurazione orbitale} $\qty(1s)^2\qty(3s)$
		
			In questa configurazione ho due elettrone in $n = 0$ ed un elettrone in $n = 2$, con $l_3 = 0$:
		
			Come prima:
			
		\begin{align*}
			E_1 = \frac{13}{2} \hbar \omega
		\end{align*}
		
			per comporre i momenti angolari distinguo in base al livello $n$, come sempre, allora:
			
		\begin{itemize}
			\item $n = 0$
				
				Per antisimmetria:
				
			\begin{align*}
				\qty(L, S)_{12} = \qty(0, 0)
			\end{align*}
			
			\item $n = 2$
			
			\begin{align*}
				\qty(L, S)_3 = \qty(0, \frac{1}{2})
			\end{align*}

		\end{itemize}

			Quindi in questa configurazione ho:
			
		\begin{align*}
			\qty(L, S) = \qty(0, 0) \otimes \qty(0, \frac{1}{2}) = \qty(0, \frac{1}{2})
		\end{align*}

			in questo modo allora abbiamo che, siccome $J = 1/2$, abbiamo degenerazione $2$.

		\item \textit{Configurazione orbitale} $\qty(1s)^2\qty(3d)$: 
		
			Ho due elettrone in $n = 0$ ed un elettrone in $n = 2$. Questa configurazione differisce dalla precedente per il valore di $l_3 = 2$. Allora:
			
		\begin{align*}
			E_1 = \frac{13}{2} \hbar \omega
		\end{align*}
		
			come al solito per comporre i momenti angoli distinguono in base al livello $n$:
		
		\begin{itemize}
			\item $n = 0$
			
				Per antisimmetria vale che:
				
			\begin{align*}
				\qty(L, S)_{12} = \qty(0, 0)
			\end{align*}

			\item $n = 2$
			
			\begin{align*}
				\qty(L, S) = \qty(0, 0) \otimes \qty(2, \frac{1}{2}) = \qty(2, \frac{1}{2})
			\end{align*}
				
		\end{itemize}
		
			In questo modo otteniamo che:
			
		\begin{align*}
			\qty(L, S) = \qty(0, 0) \otimes \qty(2, \frac{1}{2}) = \qty(2, \frac{1}{2})
		\end{align*}
		
			allora otteniamo:
				
		\begin{align*}
			J = \frac{3}{2} \oplus \frac{5}{2}
		\end{align*}
			
			con degenerazione $10$.

	\end{enumerate}

		In conclusione, la degenerazione complessiva nel primo livello eccitato è $30 + 10 + 2 = 42.$ 
	
\end{itemize}

\subsection{Sistema di $4$ Fermioni Identici nel Potenziale Armonico}
	
	Considero $4$ elettroni nel potenziale armonico:
	
\begin{align*}	
	\hat{H}_0 = \hat{H}^{\qty(1)} + \hat{H}^{\qty(2)} + \hat{H}^{\qty(3)} + \hat{H}^{\qty(4)}
\end{align*}

	con, come sappiamo:
	
\begin{align*}
	\hat{H}^{\qty(i)} = \frac{\hat{p}^2_i}{2m} + \frac{1}{2} m \omega^2 r^2_i
\end{align*}

	dove $\hat{H}^{\qty(i)}$ è l'Hamiltoniana di singola particella. Sia la base di singola particella $\qty{\ket{n, l, m, \sigma}}$.
	\\
	\\
	Analizziamo allora:
	
\begin{itemize}

	\item \textbf{Livello fondamentale} ($n_1 = n_2 = 0, n_3 = n_4 = 1$ e $l_1 = l_2 = 0, l_3 = l_4 = 1$)
	
		Al livello fondamentale ho due elettroni in $n = 0$ e due in $n = 1$. Quindi l'energia del livello sarà:
		
	\begin{align*}	
		E_0 = 8 \hbar \omega
	\end{align*}
	
		con degenerazione pari a:
		
	\begin{align*}
		\text{dim} \mathcal{H}_- = \mqty(2 \\ 2) \, \mqty(6 \\ 2) = 15
	\end{align*}
	
		il livello fondamentale ha allora degenerazione $15$.
		\\
		\\
		Componendo i momenti angolari, distinguo per livelli, ho che:
		
	\begin{itemize}
		\item $n = 0$
		
			Per antisimmetria si ottiene:
		\begin{align*}
			\qty(L, S, J)_{12} = \qty(0, 0, 0)
		\end{align*}

		\item $n = 1$
		\begin{align*}
			& L_{34} = l_3 \otimes l_4 = 1 \otimes 1 = 2 \oplus 1 \oplus 0 \\
			& S_{34} = s_3 \otimes s_4 = \frac{1}{2} \otimes \frac{1}{2} = 0 \oplus 1
		\end{align*}
		
			Allora vale che:
			
		\begin{align*}
			\qty(L, S)_{34} = \qty(0, 0) \oplus \qty(1, 1) \oplus \qty(2, 0)
		\end{align*}

	\end{itemize}

		da cui allora si ottiene che
		
	\begin{align*}
		J = 0 \oplus 0 \oplus 1 \oplus 2 \oplus 2
	\end{align*}
	
		con degenerazione $15$.

	\item \textbf{Primo livello eccitato}
	
		Al primo livello eccitato ho $3$ possibili configurazioni. Le prime $2$, differendo per un diverso valore di $l$ in $n = 2$, le faremo insieme:
		
	\begin{enumerate}
		\item \textit{Configurazioni orbitali} $\qty(1s)^2\qty(2p)\qty(3s)$ \textit{e} $\qty(1s)^2\qty(2p)\qty(3d)$
			
			In questo modo ho così due elettroni in $n = 0$, uno in $n  =1$ ed un elettrone in $n = 2$ (che può avere $l = 0$ oppure $l = 2$). In questo caso l'energia è:
			
		\begin{align*}
			E_1 = 9 \hbar \omega
		\end{align*}
		
			per quanto visto nei due esempi precedenti, possiamo già dire che in questa configurazione:
			
		\begin{align*}
			\qty(L, S) & = \qty(0, 0) \otimes \qty(1, \frac{1}{2}) \otimes \qty(\qty(0, \frac{1}{2}) \oplus \qty(2, \frac{1}{2}) ) \\
			& = \qty(1, 0) \oplus \qty(1, 1) \oplus \qty(1, 0) \oplus \qty(2, 0) \oplus \qty(3, 0) \oplus \qty(1, 1) \oplus \qty(2, 1) \oplus \qty(3, 1)
		\end{align*}
		
			da cui si ottiene:
			
		\begin{align*}
			J = 1 \oplus 0 \oplus 1 \oplus 2 \oplus 1 \oplus 2 \oplus 3 \oplus 0 \oplus 1 \oplus 2 \oplus 1 \oplus 2 \oplus 3 \oplus 2 \oplus 3 \oplus 4 \, \Rightarrow \, \text{degenerazione} = 72
		\end{align*}
	
		\item \textit{Configurazione orbitale} $\qty(1s)\qty(2p)^3$
		
			In questa configurazione abbiamo invece un elettrone in $n = 0$ e tre elettroni in $n = 1$, con rispettivamente $ l =1$ ed $m = -1, 0, +1$. 
			\\
			\\
			Anche in questo caso ho che:
			
		\begin{align*}
			E_1 = 9 \hbar \omega
		\end{align*}

			la degenerazione della configurazione è:
			
		\begin{align*}
			\text{dim} \, \mathcal{H}_- = \mqty(2 \\ 1) \, \mqty(6 \\ 3) = 2 \cdot 20 =  40
		\end{align*}

			Per comporre i momenti angolari distinguo in base al livello $n$:
			
		\begin{itemize}
			\item $n = 0$
			
				Ho un solo elettrone e quindi:
				
			\begin{align*}
				\qty(L, S)_{n = 0} = \qty(0, \frac{1}{2}) \, \Rightarrow \, J_1 = \frac{1}{2}
			\end{align*}
			
				che appunto ha degenerazione $2$.	
			
			\item $n = 1$
			
				Qui so che la generazione totale del livello è $20$, ma notiamo che:
			
			\begin{align*}
				& L_{234} = l_2 \otimes l_3 \otimes l_4 = 1 \otimes 1 \otimes 1 = 1 \oplus 0 \oplus 1 \oplus 2 \oplus 1 \oplus 1 \oplus 2 \oplus 3 \\
				& S_{234} = s_2 \otimes s_3 \otimes s_4 = \frac{1}{2} \otimes \frac{1}{2} \otimes \frac{1}{2} = \qty(0 \oplus 1) \otimes \frac{1}{2} = \frac{1}{2} \oplus \frac{1}{2} \oplus \frac{3}{2}
			\end{align*}
			
				ne segue che sicuramente una coppia antisimmetrica è:
				
			\begin{align*}
				\qty(0, \frac{3}{2})
			\end{align*}

				con degenerazione $4$. 
			\\
			\\
			Gli altri $16$ stati di questa configurazione hanno $S = \frac{1}{2}$ poiché certamente $S = \frac{3}{2}$ si esaurisce con $L = 0$, poiché totalmente simmetrico (il più alto) va con totalmente antisimmetrico (il più basso). Inoltre $L = 3$ va scartato, poiché l'unico stato totalmente simmetrico di $L$, gli altri sono misti o $0$ (unico antisimmetrico), poiché non può legarsi né allo stato totalmente simmetrico $S = 3/2$ né a quelli con simmetria mista con $S = 1/2$. In conclusione si ottiene allora:
			
		\begin{align*}
			\qty(L, S)_{n = 1} = \qty(0, \frac{3}{2}) \oplus \qty(1, \frac{1}{2}) \oplus \qty(2, \frac{1}{2})
		\end{align*}
		
			allora:
			
		\begin{align*}
			J_{n = 1} = \frac{3}{2} \oplus \frac{1}{2} \oplus \frac{3}{2} \oplus \frac{3}{2} \oplus \frac{5}{2}
		\end{align*}

		\end{itemize}
	
		In questo modo si ha che:
		
	\begin{align*}
		\qty(L, S) = \qty(0, 1) \oplus \qty(0, 2) \oplus \qty(1, 0) \oplus \qty(1, 1) \oplus \qty(2, 0) \oplus \qty(2, 1)
	\end{align*}
	
	
		da cui:
		
	\begin{align*}
		J = 1 \oplus 2 \oplus 1 \oplus 1 \oplus 0 \oplus 1 \oplus 2 \oplus 2 \oplus 1 \oplus 2 \oplus 3
	\end{align*}
		
	\end{enumerate}

\end{itemize}
	
	In conclusione, la degenerazione complessiva del primo livello eccitato è $40 + 72 = 112$.
	
	
	
	
	
	
	
	
	
	
	